%! The Statistical Cycle and Types of Data — Unit 1, Lesson 1.1 — Group Activity
\documentclass[10pt]{article}
\usepackage{saar-article}
\usepackage{saar-boxes}
\newcommand{\TallMath}[1]{$\displaystyle #1\rule[-1.4em]{0pt}{3.2em}$}

\begin{document}

\pageheader{Unit 1, Lesson 1.1}{Group Activity}
% NAMESTRIP: no \namedateperiod here. The student writes their name once, on the
% cover this component is stapled behind. See references/conventions.md.

\vspace{0.15in}

\begin{headlinebox}{frost}
\small
\textbf{Run the Cycle --- The Riverbend Breakfast Cart.} The cafeteria manager put
a breakfast cart in the main hallway for one week and wants to know whether to
keep it. Every answer needs a \emph{reason}, and every conclusion has to name the
group it applies to.
\end{headlinebox}

\vspace{0.1in}

\boxguard[30]
\begin{tcolorbox}[colback=white, colframe=black!40, breakable,
    title=\textbf{Tier R --- Remediate},
    fonttitle=\bfseries, arc=2mm, left=3mm, right=3mm, top=2mm, bottom=2mm]
\small
On Friday the manager asked $40$ students leaving the cart, ``What did you buy at
the cart this morning?''

\begin{enumerate}[leftmargin=1.6em, itemsep=6pt, topsep=4pt]

  \item Match each part of the manager's study to its stage. Write $1$, $2$, $3$,
        or $4$.

        \begin{center}
        \renewcommand{\arraystretch}{1.4}
        \begin{tabularx}{0.95\linewidth}{@{} X c @{}}
        \toprule
        \textbf{What the manager did} & \textbf{Stage} \\
        \midrule
        Asked 40 students what they bought.            & \blank{1.2cm} \\
        Decided to ask what students buy at the cart.  & \blank{1.2cm} \\
        Told the principal what the results showed.    & \blank{1.2cm} \\
        Counted the answers into a table by item.      & \blank{1.2cm} \\
        \bottomrule
        \end{tabularx}
        \end{center}

  \item Complete the frequency table. The first row is done for you.

        \begin{center}
        \renewcommand{\arraystretch}{1.4}
        \begin{tabularx}{0.9\linewidth}{@{} l c X @{}}
        \toprule
        \textbf{What they bought} & \textbf{Count} &
        \textbf{Relative frequency (percent)} \\
        \midrule
        Bagel      & 14 & $14/40 = 0.35 = 35\%$ \\
        Fruit cup  & 10 & \blank{4.2cm} \\
        Muffin     & 12 & \blank{4.2cm} \\
        Nothing    &  4 & \blank{4.2cm} \\
        \bottomrule
        \end{tabularx}
        \end{center}

  \item What was the \textbf{variable of interest} --- what did the manager
        record for each student?

        Variable: \blank{6cm} \hfill Type: \blank{3cm}

  \item Which item was bought by the most students? \blank{3.5cm}

  \item Write one sentence saying what the study found. Name the group and the
        percent.
        \writeline
\end{enumerate}
\end{tcolorbox}

\vspace{0.1in}

\boxguard
\begin{tcolorbox}[colback=white, colframe=black!40, breakable,
    title=\textbf{Tier A --- Approaching Proficiency},
    fonttitle=\bfseries, arc=2mm, left=3mm, right=3mm, top=2mm, bottom=2mm]
\small
\begin{enumerate}[leftmargin=1.6em, itemsep=5pt, topsep=3pt]

  \item The manager drafts two questions for a second study.

        \begin{itemize}[leftmargin=1.6em, itemsep=2pt, topsep=2pt]
          \item[\textbf{A.}] ``Is the breakfast cart good?''
          \item[\textbf{B.}] ``How many mornings per week do Riverbend students
                buy breakfast at school?''
        \end{itemize}

        \textbf{(a)} Which one is a statistical question? \blank{1.5cm}

        \textbf{(b)} Name one test the other question fails, and repair it.
        \writeline

        \textbf{(c)} For question \textbf{B}, what is the variable of interest,
        and what type is it?

        Variable: \blank{5cm} \hfill Type: \blank{3cm}

  \item The manager surveys $80$ Riverbend students. Of those, $24$ say they
        would buy breakfast at school every morning if the cart stayed.

        \textbf{(a)} What percent of the students surveyed is that?

        \begin{work}
          \dfrac{24}{80} &= 0.30 \\
                         &= 30\%
        \end{work}

        \textbf{(b)} Riverbend has $900$ students. Predict how many would buy
        breakfast every morning.

        \begin{work}
          0.30 \times 900 &= 270 \text{ students}
        \end{work}

        \textbf{(c)} Write the conclusion as one sentence the manager could hand
        to the principal. Name the group, the number, and what was measured.
        \writeline

  \item Each report below fails at one stage of the cycle. Write the stage number
        and say in a few words what went wrong.

        \begin{center}
        \renewcommand{\arraystretch}{1.5}
        \begin{tabularx}{0.97\linewidth}{@{} X c p{4.2cm} @{}}
        \toprule
        \textbf{Report} & \textbf{Stage} & \textbf{What went wrong} \\
        \midrule
        Only students standing at the cart were asked, and the report says
        ``Riverbend students love the cart.''
          & \blank{1cm} & \blank{3.8cm} \\
        The manager wrote down all 40 answers as a paragraph and never counted
        them by item.
          & \blank{1cm} & \blank{3.8cm} \\
        \bottomrule
        \end{tabularx}
        \end{center}
\end{enumerate}
\end{tcolorbox}

\vspace{0.1in}

\boxguard
\begin{tcolorbox}[colback=white, colframe=black!40, breakable,
    title=\textbf{Tier E --- Extension},
    fonttitle=\bfseries, arc=2mm, left=3mm, right=3mm, top=2mm, bottom=2mm]
\small
\begin{enumerate}[leftmargin=1.6em, itemsep=5pt, topsep=3pt]

  \item The school newspaper prints this:

        \begin{center}
        \textit{``$35\%$ of Riverbend students buy a bagel every morning.''}
        \end{center}

        The arithmetic behind the $35\%$ is correct. Explain \emph{two} separate
        things the sentence claims that the Tier R data cannot support.
        \writelines{2}

  \item Your group gets one hour on the district agenda. Choose a question about
        students at your school that data could settle, then plan all four
        stages.

        \begin{center}
        \renewcommand{\arraystretch}{1.35}
        \begin{tabularx}{0.97\linewidth}{@{} p{3.3cm} X @{}}
        \toprule
        Stage 1 --- our question & \blank{9cm} \\
        Variable of interest & \blank{9cm} \\
        Type of data & \blank{9cm} \\
        Stage 2 --- who we ask, and how many & \blank{9cm} \\
        Stage 3 --- how we organize it & \blank{9cm} \\
        Stage 4 --- the sentence we hope to write & \blank{9cm} \\
        \bottomrule
        \end{tabularx}
        \end{center}

  \item Two groups ask about the same students. Group 1 records \emph{how many
        mornings per week} each student buys breakfast. Group 2 records only
        whether each student \emph{ever} buys breakfast --- yes or no.

        Type for Group 1: \blank{3.2cm} \hfill Type for Group 2: \blank{3.2cm}

        The two groups collected data from the same students on the same topic.
        Explain how the \emph{question} --- not the students --- decided the type.
        \writelines{2}
\end{enumerate}
\end{tcolorbox}

\end{document}
